% Generated draft: requires mathematical and source review.
\documentclass[11pt]{article}
\usepackage[margin=25mm]{geometry}
\usepackage{fontspec}
\setmainfont{Latin Modern Roman}
\usepackage{amsmath,amssymb,mathtools,unicode-math}
\setmathfont{Latin Modern Math}
\usepackage{enumitem,needspace,xurl,hyperref}
\setlength{\parindent}{0pt}
\setlength{\parskip}{5pt}
\setlength{\emergencystretch}{3em}
\newcommand{\N}{\mathbb{N}}
\newcommand{\Z}{\mathbb{Z}}
\newcommand{\Q}{\mathbb{Q}}
\newcommand{\R}{\mathbb{R}}
\newcommand{\C}{\mathbb{C}}
\newcommand{\F}{\mathbb{F}}
\newcommand{\A}{\mathbb{A}}
\newcommand{\PP}{\mathbb{P}}
\newcommand{\E}{\mathbb{E}}
\newcommand{\eps}{\varepsilon}
\newcommand{\dd}{\,\mathrm{d}}
\begin{document}
\section*{UnsolvedMath ID: 30003637}
\textbf{Source problem code:} catalogue entry
\subsection{Definitions and mathematical statement}
Use the standard discrete Beurling interpretation: let \(\mathcal P=(p_j)_{j\geq1}\) be a sequence of labelled generalized primes satisfying \(1<p_1\leq p_2\leq\cdots\) and \(p_j\to\infty\). A generalized integer is a formal product \(n=\prod_j p_j^{a_j}\), where \(a_j\in\Z_{\geq0}\) and only finitely many \(a_j\) are nonzero. Its numerical value is \(v(n)=\prod_j p_j^{a_j}\). Let \(\mathcal N\) denote these formal products, including the empty product of value \(1\); distinct formal products are counted separately even when their numerical values coincide. Define
\[
N(x)=\#\{n\in\mathcal N:v(n)\leq x\},\qquad x\geq1,
\]
and define the generalized Möbius function by
\[
\mu_{\mathcal P}(n)=
\begin{cases}
(-1)^{\sum_j a_j},&a_j\in\{0,1\}\text{ for every }j,\\
0,&a_j\geq2\text{ for some }j.
\end{cases}
\]
Thus \(\mu_{\mathcal P}(1)=1\), and
\[
M_{\mathcal P}(x)=\sum_{\substack{n\in\mathcal N\\v(n)\leq x}}\mu_{\mathcal P}(n).
\]
Assume
\[
N(x)=x^{1+o(1)}\qquad(x\to\infty),
\]
meaning that \(\log N(x)/\log x\to1\).

\textbf{Problem.} Among Beurling generalized integer systems satisfying this counting hypothesis, determine how small the asymptotic growth of \(M_{\mathcal P}(x)\) can be.

\textbf{Unresolved formulation qualifications.} The source does not specify a comparison scale for ``how small'', such as an optimal power exponent, an eventual bound for \(\lvert M_{\mathcal P}(x)\rvert\), or an oscillation criterion. Consequently, it supplies an extremal research question rather than a fully specified conjectured inequality. The discrete, multiplicity-counted convention above makes the summation precise, but the supplied record does not explicitly state that convention. Its literature assessment mentions an exponent below one half and a proposed nonperiodic zero-set improvement without defining the associated bound, zero set, or improvement; these remarks do not determine an additional precise conjecture.
\subsection{Short English statement}
For a Beurling system whose generalized integers have counting function \(N(x)=x^{1+o(1)}\), ask how slowly its generalized Möbius partial sums can grow. The source leaves the exact measure of smallness unspecified. This statement uses the standard discrete interpretation with multiplicities and requires human mathematical review.
\subsection{Sources}
\begin{itemize}
\item Local UnsolvedMath registry: supplied complete record, permanent ID \(30003637\), title \emph{Small Möbius Sums in Generalized Integer Systems}, exact source problem code \texttt{catalogue entry}, in the Oberwolfach Reports Open Problems collection. The record gives proposed year 2017 and names no proposer. Its \texttt{open} status and \texttt{corrected\_verified} statement label are supplied registry assessments, not independent verification here.
\item Supplied report citation: \emph{Mini-Workshop: Interplay between Number Theory and Analysis for Dirichlet Series (2018)}, Oberwolfach Reports, pp.\ 3035--3069, DOI \texttt{10.4171/owr/2017/51}. \url{https://doi.org/10.4171/owr/2017/51}. The background describes Oberwolfach Reports as documenting mathematical workshops and their open-problem discussions. The supplied title includes 2018, whereas the record also identifies this as Report 51/2017 and gives proposed year 2017; that dating difference is unresolved.
\item Supplied literature source: Neamah--Hilberdink, \url{https://arxiv.org/abs/1901.06866}.
\item Supplied literature source: Vindas, \url{https://arxiv.org/abs/2406.00736}.
\item Supplied literature source: Oberwolfach Report 51/2017, \url{https://doi.org/10.4171/owr/2017/51}.
\item Literature assessment checked at \texttt{2026-08-21}: supplied dated curation, not independently verified. It labels the problem open and reports that later work proves lower bounds and special-case estimates for Beurling generalized numbers, but that no general construction with exponent below one half or the proposed nonperiodic zero-set improvement resolves the extracted problem. The background explicitly calls this dated literature triage derived from the supplied OWR curation and requests independent verification of time-sensitive status and cited claims.
\end{itemize}
\end{document}
